\documentclass{article}
\usepackage[utf8]{inputenc}
\usepackage[T1]{fontenc}
\usepackage{geometry}
\geometry{
    a4paper,
    left=2.5cm,
    right=2.5cm,
    top=2.5cm,
    bottom=2.5cm
}
\usepackage{graphicx}
\usepackage{float}
\usepackage{hyperref}
\usepackage[english]{babel}
\usepackage{amsmath}
\usepackage{amssymb}
\usepackage{array}
\usepackage{booktabs}
\usepackage{caption}
\usepackage{subcaption}
\usepackage{enumitem}
\usepackage{fancyhdr}
\usepackage{titlesec}
\usepackage{url}
\usepackage{textcomp}
\usepackage{xcolor}
\usepackage{listings}
\usepackage{multirow}
\usepackage{longtable}
\usepackage{tabularx}
\usepackage{times}

% Set up page style
\pagestyle{fancy}
\fancyhf{}
\rfoot{\thepage}
\renewcommand{\headrulewidth}{0pt}

% Title formatting
\titleformat{\chapter}[display]{\normalfont\bfseries}{\chaptertitle\thechapter}{0pt}{\Huge}
\titlespacing*{\chapter}{0pt}{50pt}{40pt}

% Document info
\title{Template Document}
\author{Academic Support Agent}
\date{\today}

\begin{document}
\maketitle

% Definitions.
% ------------
\buecedefinitions%
        {A VERY VERY VERY LONG LONG LONG TITLE TITLE TITLE TITLE OF THE REPORT}
        {AN ABBREVIATED TITLE FOR RUNNING HEADERS}
        {FullAuthor1, FullAuthor2 and FullAuthor3}
        {Date}
        {YYYY-NN} % Number of the report (four year digits and number)

% Box with title to fit the opening in the cover
% (adds an empty page in double-sided printing mode).
% ---------------------------------------------------
\buecereporttitleboxpage

% Title page
% (adds an empty page in double-sided printing mode).
% ---------------------------------------------------
\buecereporttitlepage

% Special page, e.g., if the report is restricted or
% to whom it is dedicated, etc., otherwise skip.
% (adds an empty page in double-sided printing mode).
% ---------------------------------------------------
\bueceprefacepage{Here comes a special preface page. For example, if the report
is restricted, then a suitable note can be included. This page can also be used
to indicate to whom the document is dedicated, etc.}

% Report summary; max. 1 page.
% (adds an empty page in double-sided printing mode).
% ---------------------------------------------------
\pagenumbering{roman}
\setcounter{page}{1}
\buecereportsummary{Summary of the report. Maximum 1 page.}

% Table of contents, list of figures and list of tables.
% ``\bueceemptypage'' adds empty page in double-sided
% printing mode and performs ``
% mode.
% ------------------------------------------------------




% Switch on running headers for the report:
%   odd pages  - title (lowercase); if too long, use
%                the first few words followed by ``...'',
%   even pages - last names of the authors.
% -------------------------------------------------------
\buecereportheaders

% Introduction.
% -------------
\pagenumbering{arabic}
\setcounter{page}{1}

\section{Introduction}  % Article style
%\chapter{Introduction}  % Report style

Here comes the introduction and reference to a paper \cite{Konr01cm}.

% Following sections, subsections, etc.
% -------------------------------------
\section{Starting chapter}  % Article style
%\chapter{Starting chapter}  % Report style

This is how the first chapter starts.

\subsection{Early section}  % Article style
%\section{Early section}  % Report style

And this is the first section of this chapter.

\section{Another chapter}  % Article style
%\chapter{Another chapter}  % Report style

\subsection{New section}  % Article style
%\section{New section}  % Report style

Section goes here ...

\section{Final chapter}  % Article style
%\chapter{Final chapter}  % Report style

\subsection{Another section}  % Article style
%\section{Another section}  % Report style

Section with a figure (Fig.~\ref{fig:example}).

% Plots (PostScript files) are included through the ``figure'' environment.
% For more complicated figures use the minipage commaned (see LaTeX manual).
% --------------------------------------------------------------------------
\begin{figure}[htb]
%
  \begin{minipage}[t]{0.49\linewidth}\centering
%    \centerline{\epsfig{figure=figures/regbsdcod.eps,width=8.0cm}}
    \Ovalbox{\vbox to 1.5in{\vfill\hbox{\vtop{\hsize=2.5in\hfill}\hfill}\vfill}}
    \medskip
    \centerline{(a)}
  \end{minipage}\hfill
%
  \begin{minipage}[t]{0.49\linewidth}\centering
%    \centerline{\epsfig{figure=figures/regbsdcod.eps,width=8.0cm}}
    \Ovalbox{\vbox to 1.5in{\vfill\hbox{\vtop{\hsize=2.5in\hfill}\hfill}\vfill}}
    \medskip
    \centerline{(b)}
  \end{minipage}

  \bigskip

  \begin{minipage}[t]{0.49\linewidth}\centering
%    \centerline{\epsfig{figure=figures/regbsdcod.eps,width=8.0cm}}
    \Ovalbox{\vbox to 1.5in{\vfill\hbox{\vtop{\hsize=2.5in\hfill}\hfill}\vfill}}
    \medskip
    \centerline{(c)}
  \end{minipage}\hfill
%
  \begin{minipage}[t]{0.49\linewidth}\centering
%    \centerline{\epsfig{figure=figures/regbsdcod.eps,width=8.0cm}}
    \Ovalbox{\vbox to 1.5in{\vfill\hbox{\vtop{\hsize=2.5in\hfill}\hfill}\vfill}}
    \medskip
    \centerline{(d)}
  \end{minipage}
%
  \caption{Block diagram: (a) one; (b) two; (c) three, and (d) four.}
  \label{fig:example}
\end{figure}

% Bibliography.
% -------------
\parskip=0pt
\parsep=0pt


% Important: substitute your BiBTeX (*.bib) files below.
% ------------------------------------------------------
\end{document}